\clearpage
\section*{Where aldermen differ in the zoning record}
\textit{September 29, 2026. Agent-drafted work on branch \texttt{explore-alderman-measures}; Jacob asked why the
alderman stall measures are so noisy and whether other cuts of the rezonings carry more signal.}

With the Journals (2000--2010) and eLMS (2011 to September 2026) joined
(\path{tasks/audits/zoning_stalls_and_delays}), whether an application stalls carries almost no alderman signal.
Adjusted for year, kind of change and months left in the council term, the variance of aldermen's true stall effects
is close to zero: among aldermen with at least 20 applications its standard deviation is 1 to 2 points, against 4 to 5
points of sampling error in each alderman's average, and a test that all aldermen share one rate is not rejected (p
near 0.1 in each period). Nearly every application that reaches the Council passes, consistent with applications
being screened with the alderman before they are filed. The production eLMS measure
(\path{tasks/estimate_alderman_zoning_measures}) overstates the spread because aldermen with a few applications at the
edge of the data window, often lapsing with a term, inflate its estimate of the true variance (8.5 points with them,
near zero without).

Delays carry signal, but much of it is the project. The share of applications not passed within 90 days differs
across aldermen by about 8 points beyond sampling error (p below $10^{-15}$). Large projects are slow everywhere: 84
percent of applications for 200 or more units, from the application forms
(\path{tasks/extract_zoning_application_forms}), were not passed within 90 days, against about 20 percent of small
ones. Controlling for distance to the Loop by kind of change cut Burnett's excess from 16 to 6 points and turned
Reilly's negative. Conway's applications of 2023--2026 are nearly all large downtown planned developments passed as
substitutes, slow (76 percent not passed within 90 days) but none stalled; the forms stop in mid-2023, when eLMS began
attaching them separately, so these cannot yet be adjusted for size.

What is filed and built responds more clearly to the alderman
(\path{tasks/audits/alderman_turnover_volumes}). Across the elections of 2007, 2011 and 2019, when wards kept their
boundaries, wards whose alderman changed departed from their side of the city's split of new construction permits
much more than wards whose alderman stayed (permutation p = 0.010), with falls and rises alike; applications point
the same way but not clearly (p = 0.14). Aldermen's own downzonings vary across aldermen several times more than
counting noise allows (2.15 a year in the 2000s, standard deviation 2.36 against 0.57), led by Matlak, Ocasio, Flores,
Shiller and Balcer, but they do not shift at turnovers: they come in bursts, in the same wards whoever holds them. Own
downzonings fell from 540 in the term of 2003--2007 to 46 in that of 2019--2023.

\textit{Open.} The size adjustment for 2023--2026 needs eLMS's separately attached applications, and the permit test
for those years needs permits after 2022; both are new downloads.
